\begin{quadro*}[tbp]
\caption{Desenhos de comparação e alcance da evidência obtida}\label{qua:desenhos}
\begingroup\singlespacing\fontsize{10}{12}\selectfont
\begin{tabularx}{\textwidth}{|P{.24\textwidth}|P{.34\textwidth}|X|}
\hline\textbf{Desenho} & \textbf{O que o contraste avalia} & \textbf{Limite para a interpretação}\\\hline
Sistemas completos em paralelo & Entrada, transporte, tratamento e determinação de massa, com períodos compatíveis \parencite{rsl201,rsl144}. & A diferença combina as características dos dois sistemas e suas condições de operação.\\\hline
Leitura automática e pesagem do mesmo filtro & Determinação da massa depositada, mantendo a coleta comum \parencite{rsl214}. & Não valida independentemente a coleta de dois sistemas completos.\\\hline
Subconjunto externo de estudo de exposição & Relação entre instrumentos no sítio e no período de calibração externos \parencite{rsl292,rsl427}. & Resultados internos ou posteriores já corrigidos não ampliam o número de validações.\\\hline
Rede com agregação por monitor & Diferenças entre concentrações pareadas agregadas, preservando a unidade de análise \parencite{rsl026}. & Estações e médias anuais não equivalem a pares diários ou a experimentos independentes.\\\hline
Massa inferida de tamanho de partículas & Desempenho da cadeia de medição e das hipóteses de conversão em massa \parencite{rsl288,rsl283}. & Densidade, forma e seleção por tamanho participam do resultado; incerteza não é viés percentual.\\\hline
\end{tabularx}
\Fonte{Elaboração própria a partir dos relatos citados.}
\endgroup
\end{quadro*}
